% ============================================================
% !TeX root = ../main.tex
%  Chapter 4 — Results and Evaluation
% ============================================================
\chapter{Results and Evaluation}
\label{ch:results}

\noindent This chapter presents the experimental evaluation of the approach
described in \cref{ch:methodology}.

\section{Experimental Setup}
\label{sec:setup}
\noindent
We evaluate the final CardioSDF model, obtained by pre-training on the synthetic
SSM meshes and then fine-tuning on the real end-diastolic (ED) and end-systolic
(ES) data, as described in \cref{ch:methodology}. The datasets, quality gates,
and patient-level splits are reported in \cref{subsec:dataset-composition}; all
numbers below are computed on the held-out test partitions, which contain no
patient seen during training. Two questions guide the evaluation: how faithfully
the model reconstructs the LV surfaces from sparse contours, and how the analytic
wall thickness it produces compares with independent measurement algorithms
applied to the same geometry.

\subsection{Metrics}
\label{subsec:metrics}
\noindent
Reconstruction quality is measured with four quantities. The \emph{Chamfer
distance} averages, for every reconstructed vertex, the distance to the nearest
ground-truth surface point and vice versa, reported separately for the
endocardium and the epicardium in millimetres; smaller is better. The
\emph{slice residual} is the mean absolute signed-distance value at the input
contour points, and measures how closely the reconstructed surface passes
through the observed short-axis (SAX) rings, again in millimetres. The
\emph{watertight rate} is the fraction of reconstructed meshes that form a
closed surface with no holes. The \emph{volume ratio} is the reconstructed
cavity or wall volume divided by its ground-truth value, so that a value of one
indicates no systematic over- or under-estimation.

Wall thickness is summarised by its mean, its dispersion (standard deviation),
and its 5th and 95th percentiles across the myocardial surface, all in
millimetres. These percentiles describe the thin apical and thick basal
extremes without being distorted by isolated outliers.

\subsection{Baselines}
\label{subsec:baselines}
\noindent
Because no paired ground-truth 3D mesh exists for the clinical scans (see
\cref{subsec:missing-ground-truth}), the reconstruction is compared against the
derived meshes built from the segmentations, and the analytic wall thickness is
compared against the four measurement algorithms selected in
\cref{subsec:thickness-methods}: the Laplace field method, the Yezzi--Prince
method, the SDF cone-ray method, and the Euclidean Distance Transform (EDT)
boundary sum kept as a volumetric baseline. As an external plausibility anchor,
we also report the mean wall thickness obtained directly from the input
segmentation.

\section{Results}
\label{sec:results}
\noindent
This section reports first the geometric quality of the reconstruction and then
the wall-thickness measurements obtained from the reconstructed geometry.

\subsection{Reconstruction Quality}
\label{subsec:reconstruction-quality}
\noindent
\Cref{tab:reconstruction-quality} summarises the reconstruction accuracy on the
test set. The endocardial and epicardial surfaces are recovered with a mean
Chamfer distance of $1.12$~mm and $1.17$~mm respectively, which is close to the
$1.0$~mm isotropic resolution of the prepared data and therefore near the limit
of what the derived meshes can resolve. The mean slice residual of $0.57$~mm
shows that the reconstructed surfaces pass through the observed SAX contours
with sub-millimetre agreement, confirming that the model honours its input. The
cavity and wall volumes are reproduced almost exactly, with volume ratios of
$1.00$ for the endocardium and $1.02$ for the epicardium.

Every reconstructed mesh is watertight, for both the endocardium and the
epicardium. This is a direct consequence of reading the surfaces as zero-level
sets of a continuous signed-distance field rather than assembling them from
per-slice contours, and it holds without any mesh repair, hole filling, or cap
synthesis. The monotone-epi coupling also guarantees a strictly positive wall:
the analytic thickness never collapses to zero, with a 5th percentile of
$2.33$~mm that stays well above the enforced floor.

\begin{table}[htbp]
  \centering
  \caption{Reconstruction quality of the final CardioSDF model on the test set.
    Chamfer distance and slice residual are in millimetres; the volume ratio is
    dimensionless with an ideal value of one.}
  \label{tab:reconstruction-quality}
  \begin{tabular}{l c}
    \toprule
    Metric & Value \\
    \midrule
    Endocardium Chamfer distance (mm) & 1.12 \\
    Epicardium Chamfer distance (mm)  & 1.17 \\
    Slice residual (mm)               & 0.57 \\
    Endocardium volume ratio          & 1.00 \\
    Epicardium volume ratio           & 1.02 \\
    Watertight rate (endocardium)     & 100\% \\
    Watertight rate (epicardium)      & 100\% \\
    \bottomrule
  \end{tabular}
\end{table}

\subsection{Wall-Thickness Measurement}
\label{subsec:wall-thickness-results}
\noindent
The analytic wall thickness read directly from the decoder has a mean of
$5.31 \pm 0.41$~mm across the test set, with a 5th percentile of $2.33$~mm and a
95th percentile of $6.67$~mm. These values fall within the normal range for the
LV myocardium and are consistent with the four independent algorithms applied to
the same reconstructed surfaces.

\Cref{tab:wall-thickness-results} lists the mean, standard deviation, and 5th and
95th percentiles for each of the four selected methods, and
\cref{fig:wall-thickness-comparison} shows the same comparison graphically. The
three transmural estimators agree to within roughly one millimetre of the mean:
the Laplace field reports $5.08$~mm, the SDF cone-ray method $5.47$~mm, and the
Yezzi--Prince method $4.30$~mm. The EDT boundary sum returns a lower mean of
$3.58$~mm, as expected for a voxel-based estimator that measures the shortest
transmural distance and therefore reads low in curved regions; it is retained
only as a baseline. For reference, the mean wall thickness measured directly from
the input segmentation is $3.61$~mm, and a plain nearest-neighbour measurement on
the model surface gives $4.55$~mm, so the surface-based and field-based
estimators bracket the segmentation reference from above, as expected when the
reconstruction fills in a smooth, complete wall rather than a slice-limited one.

\begin{table}[htbp]
  \centering
  \caption{Wall-thickness statistics for the four selected methods, computed on
    the CardioSDF-generated LV geometry. All values are in millimetres.}
  \label{tab:wall-thickness-results}
  \begin{tabular}{l c c c c}
    \toprule
    Method & Mean & Std. & $p_{5}$ & $p_{95}$ \\
    \midrule
    Laplace field (state of the art) & 5.08 & 1.70 & 2.92 & 8.62 \\
    Yezzi--Prince (recent)           & 4.30 & 1.43 & 0.96 & 5.80 \\
    SDF cone rays (promising)        & 5.47 & 1.15 & 3.25 & 7.16 \\
    \addlinespace
    EDT boundary sum (baseline)      & 3.58 & 0.86 & 1.49 & 4.81 \\
    \bottomrule
  \end{tabular}
\end{table}

\begin{figure}[htbp]
  \centering
  \includegraphics[width=0.9\textwidth]{wall_thickness_methods_comparison}
  \caption{Mean wall thickness of the four selected methods on the reconstructed
    LV geometry, with whiskers spanning the 5th to 95th percentile. The dashed
    line marks the mean wall thickness measured directly from the input
    segmentation.}
  \label{fig:wall-thickness-comparison}
\end{figure}

\section{Discussion}
\label{sec:discussion}
\noindent
The results support the two central design choices of CardioSDF. First, encoding
the LV as a signed-distance field yields watertight surfaces in every test case
without any post-processing, which removes the mesh-repair heuristics that a
contour-stacking pipeline would otherwise require. Second, the monotone-epi
coupling delivers a wall that is positive everywhere, so the analytic thickness
is usable as a measurement rather than merely a by-product: its mean of
$5.31$~mm and its percentile range agree with the established Laplace field and
with the SDF cone-ray estimator to within about one millimetre.

The spread between methods is informative rather than contradictory. Field-based
transmural estimators and surface-based ray estimators measure thickness along
slightly different paths, and the EDT baseline reads systematically lower because
straight-line voxel distances underestimate thickness where the wall curves. The
agreement of the three transmural methods, together with the sub-millimetre slice
residual and near-unit volume ratios, indicates that the reconstructed geometry
is both accurate at the observed slices and plausible between them, which is
precisely where sparse SAX acquisitions provide no direct evidence.

\section{Threats to Validity}
\label{sec:threats}
\noindent
The most important limitation concerns the ground truth. The reconstruction is
evaluated against meshes \emph{derived} from the segmentations, not against
independently acquired 3D surfaces, so the Chamfer distances measure agreement
with a derived target rather than with true anatomy. For the same reason, the
wall-thickness comparison is an inter-method and reproducibility study on shared
geometry, not a clinical validation: without expert manual wall-thickness
measurements, the thesis can claim an objective and locally resolved
computational framework, but not the replacement of clinical measurement
protocols.

Several further factors bound the external validity of the results. The real
data come from two public collections (ACDC and M\&Ms-2), and the ES test
partition is small, so per-phase estimates carry more uncertainty than the ED
ones. The wall-thickness percentiles depend on the AHA orientation and on the
$96^3$ inference grid, both of which introduce discretisation effects that a
finer grid or an alternative anatomical frame could shift. Finally, the synthetic
epicardium used during pre-training is a geometric construction rather than a
measured surface, which the fine-tuning stage on real data is designed to
correct but cannot fully rule out as a residual prior.
